\section{Билет 57. Идеальный газ в поле сил тяжести. Барометрическая формула с выводом. Изменение плотности газа с высотой.}

\begin{definition}
Рассмотрим идеальный газ, находящийся в однородном поле тяжести (например, атмосфера Земли). В состоянии термодинамического равновесия газ распределён по высоте неравномерно: нижние слои сжаты весом вышележащих слоёв, поэтому давление и концентрация молекул убывают с высотой. Для вывода зависимости давления от высоты рассмотрим элементарный горизонтальный слой газа толщиной $dh$ и площадью $S=1$.
\end{definition}

\begin{theorem}[Вывод барометрической формулы]
Давление на высоте $h$ равно $p$, а на высоте $h+dh$ --- $p+dp$ (где $dp<0$). Разность давлений равна весу газа в выделенном слое:
\begin{equation}
    p - (p + dp) = \rho g \, dh \quad \Rightarrow \quad -dp = \rho g \, dh,
\end{equation}
где $\rho$ --- плотность газа, $g$ --- ускорение свободного падения.
\par\noindent Используя уравнение состояния идеального газа $pV = \frac{m}{M}RT$ и определение плотности $\rho = m/V$, выразим $\rho$ через давление:
\begin{equation}
    \rho = \frac{pM}{RT},
\end{equation}
где $M$ --- молярная масса газа, $R$ --- универсальная газовая постоянная, $T$ --- абсолютная температура (считаем изотермическим, $T=\text{const}$).
\par\noindent Подставляя $\rho$ в гидростатическое уравнение, получаем:
\begin{equation}
    dp = -\frac{pMg}{RT} \, dh \quad \Rightarrow \quad \frac{dp}{p} = -\frac{Mg}{RT} \, dh. \label{eq:diff_barometric}
\end{equation}
Интегрируем от уровня $h=0$ (давление $p_0$) до произвольной высоты $h$ (давление $p$):
\begin{equation}
    \int_{p_0}^{p} \frac{dp}{p} = -\frac{Mg}{RT} \int_{0}^{h} dh \quad \Rightarrow \quad \ln\frac{p}{p_0} = -\frac{Mgh}{RT}.
\end{equation}
Потенцируя, получаем \textbf{барометрическую формулу}:
\begin{equation}
    p(h) = p_0 \exp\left(-\frac{Mgh}{RT}\right). \label{eq:barometric_p}
\end{equation}
\end{theorem}

\begin{property}[Изменение плотности и концентрации с высотой]
Так как при постоянной температуре плотность пропорциональна давлению ($\rho = pM/RT$), закон изменения плотности с высотой аналогичен:
\begin{equation}
    \rho(h) = \rho_0 \exp\left(-\frac{Mgh}{RT}\right).
\end{equation}
В молекулярно-кинетических терминах удобно использовать концентрацию молекул $n = N/V$. Учитывая, что $p = nkT$ и $R = kN_A$, $M = mN_A$, барометрическую формулу можно записать как:
\begin{equation}
    n(h) = n_0 \exp\left(-\frac{mgh}{kT}\right), \label{eq:barometric_n}
\end{equation}
где $m$ --- масса одной молекулы, $k$ --- постоянная Больцмана, $n_0$ --- концентрация на уровне $h=0$.
\end{property}

\begin{remark}
\textbf{Физическая интерпретация:}
\begin{enumerate}
    \item Давление и концентрация убывают с высотой тем быстрее, чем больше молярная масса газа $M$ (тяжёлые газы ``оседают'' ближе к поверхности).
    \item С ростом температуры $T$ экспонента становится ``пологее'': тепловое движение эффективнее противостоит гравитации, и газ распределяется более равномерно.
    \item Формула справедлива для изотермической атмосферы. В реальной атмосфере температура меняется с высотой, поэтому зависимость осложняется, но экспоненциальный характер распределения сохраняется в первом приближении.
\end{enumerate}
\end{remark}
